\indent ಅವ್ಯಯಗಳಲ್ಲಿ ಎಂಟು ವಿಭೇದಗಳುಂಟು ; ಅವುಗಳನ್ನು ವಿವರವಾಗಿ ತಿಳಿಯೋಣ.\par

\vspace{4pt plus 4pt minus 1pt}\noindent{\kbvsecfont\fontsize{11pt}{14.5pt}\selectfont\color{kbvaccentdark} (೧)\hspace{0.45em}ಸಾಮಾನ್ಯಾವ್ಯಯ :}\hspace{0.5em}ಮೆಲ್ಲಗೆ, ಬೇಗನೆ ವೇಗವಾಗಿ, ಕೂಡಲೆ ಸುಮ್ಮನೆ, ಚೆನ್ನಾಗಿ ತಟ್ಟನೆ, ಎತ್ತರವಾಗಿ, ಇಂತು ಇತ್ಯಾದಿ ಪದಗಳು ಕ್ರಿಯೆ ನಡೆಯುವ ಯನ್ನು ಹೇಳುವ ಅವ್ಯಯಗಳು. ಈ ಕ್ರಿಯಾವಿಶೇಷಗಳಿಗೆ ಸಾಮಾನ್ಯಾವ್ಯಯಗಳೆಂದು ಹೆಸರು.\par

\vspace{4pt plus 4pt minus 1pt}\noindent{\kbvsecfont\fontsize{11pt}{14.5pt}\selectfont\color{kbvaccentdark} (೨)\hspace{0.45em}ಅನುಕರಣಾವ್ಯಯ :}\hspace{0.5em}ಚಟಚಟ, ಪಟಪಟ, ತಟತಟ, ದಡದಡ, ಕರಕರ, ಸರಸರ, ದಬದಬ, ಚುರುಚುರು, ಧಗಧಗ, ರೊಯ್ಯನೆ, ಭರ್ರನೆ, ಗಿರ್ರನೆ ಸುಯ ಇವೇ ಮೊದಲಾದ ಪದಗಳು ಅರ್ಥರಹಿತ ಧ್ವನಿ ವಿಶೇಷಗಳು. ಕೇಳಿದ ರೀತಿಯಲ್ಲಿ ಅನುಕರಣ ಮಾಡಿ ಹೇಳುವ ಈ ಶಬ್ದಗಳಿಗೆ ಅನುಕರಣಾವ್ಯಯಗಳೆಂದು ಹೆಸರು.\par

\vspace{4pt plus 4pt minus 1pt}\noindent{\kbvsecfont\fontsize{11pt}{14.5pt}\selectfont\color{kbvaccentdark} (೩)\hspace{0.45em}ಭಾವಸೂಚಕಾವ್ಯಯ :}\hspace{0.5em}ಕೋಪ, ದುಃಖ, ಆಶ್ಚರ್ಯ, ತಿರಸ್ಕಾರ, ಸಂತೋಷ, ಮೆಚ್ಚು ಇತ್ಯಾದಿ ಭಾವಗಳನ್ನು ಪ್ರದರ್ಶಿಸುವಾಗ ಎಲೆಲಾ! ಅಯ್ಯೋ! ಅಬ್ಬಾ! ಥೂ! ಓಹೋ! ಭಲೇ!... ಇವೇ ಮೊದಲಾದ ಅರ್ಥರಹಿತವಾದ ಶಬ್ದಗಳನ್ನು ಪ್ರಕಟಿಸುತ್ತೇವೆ. ಈ ರೀತಿಯ ಶಬ್ದಗಳಿಗೆ ಭಾವಸೂಚಕಾವ್ಯಯಗಳೆಂದು ಹೆಸರು.\par

\kbvsechead{(೪)\hspace{0.45em}ಕ್ರಿಯಾರ್ಥಕಾವ್ಯಯ :}

\kbvsubtitle{ಕ್ರಿಯಾಪದ ಪ್ರಕರಣ ನೋಡಿ}

\vspace{4pt plus 4pt minus 1pt}\noindent{\kbvsecfont\fontsize{11pt}{14.5pt}\selectfont\color{kbvaccentdark} (೫)\hspace{0.45em}ಸಂಬಂಧಾರ್ಥಕಾವ್ಯಯ :}\hspace{0.5em}ಎರಡು ಪದಗಳನ್ನು ಜೋಡಿಸಲು “ಮತ್ತು, ಅಥವಾ, ಊ” ಎಂದು.\par
